\documentclass[11pt]{article}
\usepackage[a4paper,margin=1in]{geometry}
\usepackage{amsmath,amssymb,mathtools}
\usepackage{siunitx}
\usepackage{booktabs}
\usepackage{hyperref}
\usepackage{enumitem}

\title{THzSim2 Measured-Reference Drude Study Notes}
\author{THzSim2 Repository}
\date{\today}

\begin{document}
\maketitle

\section{Purpose}
This note documents the equations used by the measured-reference compatibility workflow in
\texttt{thzsim2}. The intended use is to explain the code path behind the one-layer Drude
study that starts from a measured reference CSV, resamples it to a uniform time grid, fits a
single Drude layer, and exports both the old-compatible mean-squared-error summary and the new
diagnostic metrics.

\section{Reference Trace Import and Resampling}
The canonical time-trace CSV schema is
\begin{center}
\texttt{time\_ps,trace}
\end{center}
but the importer also accepts common measured-data headers such as
\begin{center}
\texttt{Time\_abs/ps, Signal/nA}.
\end{center}

If the measured time axis is not exactly uniform, the workflow preserves the raw amplitude
values but interpolates them onto a uniform grid. Let the imported samples be
\[
\{(t_i, E_i)\}_{i=0}^{N-1},
\qquad
t_{i+1} > t_i.
\]
The dominant or median time step is used as the target spacing
\[
\Delta t_{\mathrm{resamp}} = \operatorname{median}(t_{i+1}-t_i),
\]
and the uniform grid is
\[
\tilde{t}_i = t_0 + i \, \Delta t_{\mathrm{resamp}},
\qquad i = 0,\dots,N-1.
\]
The resampled trace is then
\[
\tilde{E}(\tilde{t}_i) = \operatorname{interp}\!\left(\tilde{t}_i; \{t_j,E_j\}\right),
\]
using linear interpolation.

The compatibility workflow keeps the imported amplitude in source units. In the example
measured CSV this means the exported reference trace remains in \si{\nano\ampere} rather than
being normalized.

\section{Fourier Transform Convention}
The reference and sample traces are propagated through the frequency domain using the discrete
Fourier-transform convention implemented in \texttt{thzsim2.core.fft}. If the time trace is
\[
E_n = E(t_n), \qquad t_n = t_0 + n \Delta t,
\]
then the angular-frequency grid is
\[
\omega_k = 2\pi f_k,
\]
with the standard FFT ordering from the uniform grid. The forward transform in code is
\[
\hat{E}(\omega_k) = \Delta t \, N \, \operatorname{IFFT}(E_n) \, e^{i \omega_k t_0},
\]
and the inverse transform is
\[
E_n = \frac{1}{N \Delta t}\operatorname{FFT}\!\left(\hat{E}(\omega_k)e^{-i\omega_k t_0}\right).
\]

This convention preserves the original time-origin offset \(t_0\), which matters when the
reference trace is imported from a measured file rather than generated symmetrically around
zero.

\section{Single-Layer Transmission Model}
For normal incidence, the complex field transmission through one layer of thickness \(d\),
refractive index \(n_1(\omega)\), input medium \(n_0\), and output medium \(n_2\) is
\[
T(\omega)
= t_{01}(\omega)\,p_1(\omega)\,t_{12}(\omega)\,F_{\mathrm{FP}}(\omega),
\]
where the Fresnel coefficients are
\[
t_{jk} = \frac{2n_j}{n_j+n_k},
\qquad
r_{jk} = \frac{n_j-n_k}{n_j+n_k},
\]
the propagation factor is
\[
p_1(\omega) = \exp\!\left(i \omega n_1(\omega)\frac{d}{c_0}\right),
\]
and the finite Fabry--Perot correction used by the code is
\[
F_{\mathrm{FP}}(\omega)
= \sum_{m=0}^{M} \left[r_{12}(\omega)p_1(\omega)r_{10}(\omega)p_1(\omega)\right]^m.
\]
Here \(M\) is the configured maximum number of internal reflections. In the old-compatible
study the default is \(M=8\).

\section{Drude Material Model}
The single fitted layer uses the Drude dielectric function
\[
\varepsilon(\omega)
=
\varepsilon_\infty
- \frac{\omega_p^2}{\omega(\omega+i\gamma)}.
\]
The complex refractive index is obtained from
\[
\tilde{n}(\omega) = n(\omega) + i k(\omega) = \sqrt{\varepsilon(\omega)},
\]
with the passive branch selected so that
\[
k(\omega) \ge 0.
\]

The public notebook API uses \si{\tera\hertz} for both the Drude plasma frequency and damping:
\[
\omega_p = 2\pi f_p \times 10^{12},
\qquad
\gamma = 2\pi f_\gamma \times 10^{12}.
\]

\subsection{Conversion from Conductivity and Relaxation Time}
The compatibility study follows the old script and often describes true values in terms of the
DC conductivity \(\sigma\) and relaxation time \(\tau\). The code converts these to Drude
parameters by
\[
\gamma = \frac{1}{\tau},
\]
\[
\sigma = \varepsilon_0 \frac{\omega_p^2}{\gamma},
\]
so that
\[
\omega_p = \sqrt{\frac{\sigma}{\varepsilon_0 \tau}}.
\]
The notebook-facing helper functions therefore implement
\[
f_\gamma = \frac{1}{2\pi \tau},
\qquad
f_p = \frac{1}{2\pi \times 10^{12}}
\sqrt{\frac{\sigma}{\varepsilon_0 \tau}}.
\]

\section{Why the Exported \(n,k\) CSV Starts Above Zero Frequency}
For a Drude medium,
\[
\varepsilon(\omega)
=
\varepsilon_\infty
- \frac{\omega_p^2}{\omega(\omega+i\gamma)}.
\]
As \(\omega \to 0\), the conductivity-driven response becomes singular, so the apparent
\(n(\omega)\) and \(k(\omega)\) can become very large. The code still handles the
\(\omega \approx 0\) behavior internally through a protected frequency floor, but the exported
sample-layer CSVs and plots omit the exact \SI{0}{\tera\hertz} row so that the saved
\texttt{freq\_thz,n,k} files remain numerically readable and visually useful.

\section{Noise Model}
The compatibility study adds white Gaussian noise to the clean simulated sample trace. If the
requested dynamic range is \(D\) dB and the peak absolute sample amplitude is
\[
A_{\max} = \max_t |E_{\mathrm{sample,clean}}(t)|,
\]
then the noise standard deviation is
\[
\sigma_{\mathrm{noise}} = \frac{A_{\max}}{10^{D/20}}.
\]
The observed synthetic sample trace is
\[
E_{\mathrm{obs}}(t) = E_{\mathrm{sample,clean}}(t) + \eta(t),
\qquad
\eta(t) \sim \mathcal{N}(0,\sigma_{\mathrm{noise}}^2).
\]

\section{Objective Functions and Diagnostic Metrics}
\subsection{Raw Mean-Squared Error}
For compatibility with the old study, the primary fit objective remains the raw mean-squared
error
\[
\mathrm{MSE}
=
\frac{1}{N}\sum_{i=0}^{N-1}
\left(E_{\mathrm{model}}(t_i)-E_{\mathrm{obs}}(t_i)\right)^2.
\]
This value can be numerically small when the observed trace amplitude itself is small, so it
must be interpreted together with the trace plots and the scale-aware diagnostics below.

\subsection{Relative \texorpdfstring{$L^2$}{L2} Error}
The relative \(L^2\) metric exported by the workflow is
\[
\mathrm{relative\_}L^2
=
\frac{\lVert E_{\mathrm{model}}-E_{\mathrm{obs}}\rVert_2}
{\max\!\left(\lVert E_{\mathrm{obs}}\rVert_2,10^{-30}\right)}.
\]

\subsection{Windowed MSE}
The added diagnostic \texttt{windowed\_mse} focuses on the part of the waveform where the
observed signal is not negligible. With a threshold fraction \(\alpha\) (currently
\(\alpha=0.05\)),
\[
\mathcal{I}
=
\left\{i:\, |E_{\mathrm{obs}}(t_i)| \ge \alpha \max_j |E_{\mathrm{obs}}(t_j)|\right\}.
\]
Then
\[
\mathrm{windowed\_MSE}
=
\frac{1}{|\mathcal{I}|}
\sum_{i\in\mathcal{I}}
\left(E_{\mathrm{model}}(t_i)-E_{\mathrm{obs}}(t_i)\right)^2.
\]
If no samples pass the threshold, the code falls back to the raw MSE.

\subsection{Peak-Normalized RMSE}
To make the residual scale easier to interpret visually, the workflow also exports
\[
\mathrm{peak\_normalized\_RMSE}
=
\frac{\sqrt{\frac{1}{N}\sum_i \left(E_{\mathrm{model}}(t_i)-E_{\mathrm{obs}}(t_i)\right)^2}}
{\max\!\left(\max_i |E_{\mathrm{obs}}(t_i)|,10^{-30}\right)}.
\]

\section{Optimization Strategy}
The fitting code uses a two-stage optimization strategy:
\begin{enumerate}[leftmargin=2em]
\item a global differential-evolution search over the bounded fit variables;
\item a local \texttt{L-BFGS-B} refinement starting from the global best point.
\end{enumerate}
For the old-compatible study, the bounded fit variables are:
\begin{center}
\begin{tabular}{llll}
\toprule
Parameter & Initial value & Lower bound & Upper bound \\
\midrule
Thickness & \SI{150}{\micro\meter} & \SI{40}{\micro\meter} & \SI{320}{\micro\meter} \\
\(\tau\) & \SI{3}{\pico\second} & \SI{0.05}{\pico\second} & \SI{30}{\pico\second} \\
\(\sigma\) & \SI{0.01}{\siemens\per\meter} & \SI{5e-4}{\siemens\per\meter} & \SI{0.2}{\siemens\per\meter} \\
\bottomrule
\end{tabular}
\end{center}

Internally, the optimizer works on thickness, Drude plasma frequency, and Drude damping
frequency, because those are the native variables of the forward model.

\section{Parameter-Uncertainty Estimate}
After the optimum \(\hat{\mathbf{x}}\) is found, the code estimates local parameter
uncertainties with a finite-difference Jacobian of the residual vector
\[
\mathbf{r}(\mathbf{x}) = E_{\mathrm{obs}} - E_{\mathrm{model}}(\mathbf{x}).
\]
The Jacobian is approximated column-by-column as
\[
J_{:,j}
\approx
\frac{\mathbf{r}(\hat{\mathbf{x}} + h_j \mathbf{e}_j)
-
\mathbf{r}(\hat{\mathbf{x}} - h_j \mathbf{e}_j)}
{2h_j},
\]
with a bounded step size chosen from the parameter magnitude and the parameter span.

If there are \(m\) residual samples and \(p\) fitted parameters, the residual variance estimate is
\[
\hat{\sigma}^2 = \frac{\mathbf{r}^\mathsf{T}\mathbf{r}}{\max(m-p,1)}.
\]
The covariance estimate is then
\[
\Sigma \approx \hat{\sigma}^2 (J^\mathsf{T}J)^+,
\]
where \((\cdot)^+\) denotes the pseudoinverse. The reported parameter standard deviations are
\[
\sigma_j = \sqrt{\Sigma_{jj}},
\]
and the reported parameter correlations are
\[
\rho_{ij} = \frac{\Sigma_{ij}}{\sigma_i \sigma_j}.
\]

\section{Progress, ETA, and Study Size}
The old-compatible workflow evaluates a pilot set of cases spread across the full sweep,
estimates the average runtime per case, and prints an initial total runtime estimate before the
full study loop begins. For the requested sweep
\[
4 \times 20 \times 20 \times 3 = 4800
\]
cases are run, with one replicate per case.

\section{Known Pitfalls}
\begin{itemize}[leftmargin=2em]
\item \textbf{Pulse placement matters.} A generated pulse can look artificially tiny if its center
lies outside the generated time window. The notebook-facing pulse generator now rejects that
configuration explicitly.
\item \textbf{Measured traces can carry baseline and DC offset.} The compatibility workflow keeps
the measured amplitude scale and baseline for apples-to-apples comparison, but it records basic
baseline diagnostics in the reference manifest.
\item \textbf{Raw MSE is scale-dependent.} A visibly poor fit can still have a small raw MSE if the
trace amplitude is small. The workflow therefore exports \texttt{windowed\_mse},
\texttt{relative\_l2}, and \texttt{peak\_normalized\_rmse} alongside the raw MSE.
\item \textbf{Drude DC export is not plot-friendly.} The internal model still handles the low-frequency
limit, but the exported \texttt{freq\_thz,n,k} files start at the first positive frequency to avoid
an unreadable DC row.
\end{itemize}

\end{document}
